\documentclass{article}

\usepackage{../paper/neurips_2026}

\usepackage[utf8]{inputenc}
\usepackage[T1]{fontenc}
\usepackage{hyperref}
\usepackage{url}
\usepackage{booktabs}
\usepackage{amsfonts}
\usepackage{nicefrac}
\usepackage{microtype}
\usepackage{xcolor}
\usepackage{amsmath,amssymb,amsthm,mathtools}
\usepackage{bm}
\usepackage{enumitem}
\usepackage[nameinlink,noabbrev]{cleveref}

\newtheorem{theorem}{Theorem}
\newtheorem{proposition}[theorem]{Proposition}
\newtheorem{corollary}[theorem]{Corollary}
\newtheorem{lemma}[theorem]{Lemma}
\theoremstyle{definition}
\newtheorem{definition}[theorem]{Definition}
\newtheorem{remark}[theorem]{Remark}
\newtheorem{assumption}[theorem]{Assumption}

\newcommand{\R}{\mathbb{R}}
\newcommand{\E}{\mathbb{E}}
\newcommand{\PP}{\mathcal{P}}
\newcommand{\PPtwo}{\mathcal{P}_2}
\newcommand{\W}{W_1}
\newcommand{\dd}{\,\mathrm{d}}
\newcommand{\riskhead}{\mathcal{R}_{\mathrm{head}}}
\newcommand{\riskmean}{\mathcal{R}_{\mathrm{mean}}}
\newcommand{\diag}{\operatorname{diag}}
\newcommand{\Sym}{\operatorname{Sym}}

\title{A Head Mean-Field Limit for Efficient Ensembles as $K\to\infty$}

\author{Anonymous Authors}

\begin{document}

\maketitle

\begin{abstract}
We study a different asymptotic regime for efficient ensembles: instead of sending the hidden width to infinity at fixed ensemble size, we fix the width $m$ and let the number of heads $K$ grow. Motivated by a two-layer BatchEnsemble-like proxy with independent output heads and head-wise averaged loss, we identify the empirical measure of head-specific states as the natural large-$K$ object. The first key point is structural: parameter sharing does not obstruct a head mean-field description, because the pair consisting of the shared backbone state and the empirical head law is already an exact closed state variable at finite $K$. After normalizing time so that the shared dynamics remain order one, equal learning rates would make the head dynamics degenerate as $K\to\infty$; the nontrivial regime instead keeps an order-one effective head-learning scale $\lambda\in[0,\infty)$. This yields a coupled ODE-transport system for the shared parameters and the head law. The practically relevant observable is then the induced predictor law, obtained by pushing the head law through the map $z\mapsto f(\cdot;q_t,z)$: state-level non-collapse alone does not imply functional diversity, but any exact function-level collapse must come from degeneracy of that map rather than from the transport equation crushing the head law. For the concrete two-layer proxy we prove a finite-horizon large-$K$ theorem under bounded-support initialization. This note is intended as a clean starting point for a large-ensemble theory parallel to, rather than replacing, the fixed-$K$ width-mean-field analysis.
\end{abstract}

\section{Introduction}

Efficient ensembles such as BatchEnsemble and TabM share most of their parameters across ensemble members while keeping a smaller head-specific component \citep{batchensemble,tabm}. The fixed-$K$, width-to-infinity regime studied in the companion draft emphasizes one phenomenon: in a deterministic feature-learning mean-field limit with symmetric initialization, headwise predictors can collapse to a common function. That viewpoint is useful, but it is not the only asymptotic question one can ask.

This note explores a different limit: what happens when the number of heads itself grows? For the objective
\[
  \frac1K\sum_{\alpha=1}^K \ell(f_\alpha,y),
\]
the answer is not automatic. Increasing $K$ changes both the averaging seen by the shared backbone and the raw gradient scale seen by each head. In particular, with equal learning rates across shared and head-specific parameters, each head receives a gradient of order $1/K$, so the naive large-$K$ limit freezes the head states. In the present note we therefore focus on the nontrivial regime in which the head-side learning rate is scaled so that the \emph{effective} head learning remains order one as $K$ grows. In that regime the empirical measure of head states evolves by a transport equation coupled to the shared backbone. This replaces the question ``do all heads become identical?'' by the question ``what head population law is selected, and what predictor law does it induce?'' It is a different notion of collectivity from the fixed-$K$ width-mean-field collapse theorem, and it may be closer to the operational question practitioners ask when they scale the ensemble size.

This perspective is complementary to the infinite-width GP/NTK literature. In lazy or NTK-type limits, randomness at initialization can survive at order one and the training dynamics is described by a kernel or linearized model \citep{embedded-ensembles,chizat2019lazy,mei2019}. Here the width $m$ is fixed, the shared backbone can keep feature-learning, and the macroscopic object is instead a transport equation for head states evolving in a finite-dimensional common environment. The relevant ``large object'' is therefore not an initialization kernel but a time-dependent head law.

\paragraph{Main message.}
For head-wise averaged losses, taking $K\to\infty$ still fits a mean-field template, but the correct mean-field object is the empirical distribution of head states rather than a single common predictor. Parameter sharing does not break this description: it enters only through a common backbone state driven by the empirical head law. The practically relevant observable is the induced predictor law
\[
  \bigl(z\mapsto f(\cdot;q_t,z)\bigr)_{\#}\nu_t.
\]
A dispersed head law does not by itself imply functional diversity, and a Dirac predictor law can arise even when $\nu_t$ is not a Dirac mass. Under order-one head learning one obtains a nontrivial ODE-PDE system; under bounded-Lipschitz characteristic flow, exact state collapse from dispersed initial data is impossible, so any exact function-level collapse must come from degeneracy of the head-to-predictor map rather than from the transport equation itself.

\section{Setup}
\label{sec:setup}

Let $(\mathcal X\times\mathcal Y,\mathsf P)$ be the data space and fix a hidden width $m\in\{1,2,\dots\}$. We consider the two-layer BatchEnsemble-like proxy
\begin{equation}
  f_\alpha^m(x)
  =
  \frac1m\sum_{j=1}^m
  a_{\alpha,j}\,
  \sigma\bigl(w_j^\top(r_\alpha\odot x)+b_j\bigr),
  \qquad
  \alpha=1,\dots,K,
  \label{eq:model}
\end{equation}
with shared hidden parameters $(w_j,b_j)_{j=1}^m$ and head-specific parameters $((a_{\alpha,j})_{j=1}^m,r_\alpha)$.

The population objective is the head-wise averaged risk
\begin{equation}
  \mathcal R_K
  =
  \E_{(x,y)\sim\mathsf P}
  \Bigl[
    \frac1K\sum_{\alpha=1}^K \ell(f_\alpha^m(x),y)
  \Bigr].
  \label{eq:riskK}
\end{equation}
This is \emph{not} the loss of the averaged predictor. Instead, it is the expectation of the loss of a uniformly sampled head.

It is convenient to collect the shared state and the head state as
\[
  q:=((w_j,b_j)_{j=1}^m)\in\mathcal Q:=\bigl(\R^d\times\R\bigr)^m,
  \qquad
  z_\alpha:=((a_{\alpha,j})_{j=1}^m,r_\alpha)\in\mathcal Z:=\R^m\times\R^d.
\]
We then write the output of a generic head state $z=(a,r)$ against a shared state $q=(w,b)$ as
\begin{equation}
  f(x;q,z)
  :=
  \frac1m\sum_{j=1}^m a_j\,\sigma\bigl(w_j^\top(r\odot x)+b_j\bigr).
  \label{eq:generic-output}
\end{equation}

Define the shared and head vector fields
\begin{align}
  \mathcal F(q,z)
  &:=
  -\E_{(x,y)\sim\mathsf P}
  \bigl[
    \partial_1\ell(f(x;q,z),y)\,\nabla_q f(x;q,z)
  \bigr],
  \label{eq:Fdef}
  \\
  \mathcal G(q,z)
  &:=
  -\E_{(x,y)\sim\mathsf P}
  \bigl[
    \partial_1\ell(f(x;q,z),y)\,\nabla_z f(x;q,z)
  \bigr].
  \label{eq:Gdef}
\end{align}
Under bounded inputs, $\sigma\in C_b^1(\R)$, and bounded Lipschitz $\partial_1\ell(\cdot,y)$, the formulas above define continuous vector fields that are locally Lipschitz on bounded subsets of $\mathcal Q\times\mathcal Z$. However, for the literal proxy \eqref{eq:model} they are generally \emph{not} globally bounded, because the drifts grow linearly in the head coefficients $a$ and the shared weights $w$. Accordingly, the global large-$K$ results below are stated for abstract vector fields $(\mathcal F,\mathcal G)$ satisfying explicit bounded-Lipschitz assumptions. The concrete two-layer proxy re-enters later through local affine/quadratic specializations and would require either clipping or a priori bounded-region control to justify the same global hypotheses rigorously.

\subsection{Why Parameter Sharing Does Not Obstruct a Head Mean Field}

The first conceptual issue is whether one can really speak about a mean-field limit over heads when the backbone parameters are shared rather than replicated. In the present architecture, the answer is yes.

Let $\lambda_K\ge 0$ denote the effective head-learning scale. In this note we keep the head dynamics on an order-one time scale, with the canonical choice $\lambda_K\equiv 1$. More generally, we allow any sequence $\lambda_K\to\lambda\in[0,\infty)$. Consider the finite-$K$ dynamics
\begin{align}
  \dot q_t^K
  &=
  \frac1K\sum_{\alpha=1}^K \mathcal F(q_t^K,z_{\alpha,t}^K),
  \label{eq:general-q}
  \\
  \dot z_{\alpha,t}^K
  &=
  \lambda_K\,\mathcal G(q_t^K,z_{\alpha,t}^K),
  \qquad
  \alpha=1,\dots,K.
  \label{eq:general-z}
\end{align}
Define the empirical measure of head states
\[
  \nu_t^K := \frac1K\sum_{\alpha=1}^K \delta_{z_{\alpha,t}^K}.
\]

\begin{proposition}[Exact empirical closure despite parameter sharing]
\label{prop:closure}
For every fixed $K$ and every $\lambda_K\ge 0$, the pair $(q_t^K,\nu_t^K)$ is an exact closed state variable for \eqref{eq:general-q}--\eqref{eq:general-z}. More precisely,
\begin{align}
  \dot q_t^K
  &=
  \int_{\mathcal Z}\mathcal F(q_t^K,z)\,\nu_t^K(\dd z),
  \label{eq:closure-q}
  \\
  \partial_t \nu_t^K
  +
  \lambda_K\nabla_z\cdot\bigl(\nu_t^K\,\mathcal G(q_t^K,z)\bigr)
  &=
  0
  \label{eq:closure-nu}
\end{align}
in the sense of distributions. In particular, whenever the labeled ODE system \eqref{eq:general-q}--\eqref{eq:general-z} is uniquely solvable on a time interval, two labeled head configurations with the same initial shared state $q_0^K$ and the same empirical measure $\nu_0^K$ generate the same trajectories $q_t^K$ and $\nu_t^K$ on that interval.
\end{proposition}

\begin{proof}
The identity \eqref{eq:closure-q} is just the definition of $\nu_t^K$. For the transport equation, test against a smooth compactly supported function $\varphi$:
\[
  \frac{\dd}{\dd t}\int_{\mathcal Z}\varphi(z)\,\nu_t^K(\dd z)
  =
  \frac1K\sum_{\alpha=1}^K
  \nabla\varphi(z_{\alpha,t}^K)\cdot \lambda_K\mathcal G(q_t^K,z_{\alpha,t}^K)
  =
  \int_{\mathcal Z}
  \nabla\varphi(z)\cdot \lambda_K\mathcal G(q_t^K,z)\,\nu_t^K(\dd z),
\]
which is the weak form of \eqref{eq:closure-nu}. For the last claim, note that \cref{eq:general-q,eq:general-z} is permutation-invariant in the head labels: if $\pi$ is a permutation of $\{1,\dots,K\}$ and
\[
  \widehat z_{\alpha,t}^K:=z_{\pi(\alpha),t}^K,
\]
then $(q_t^K,(\widehat z_{\alpha,t}^K)_{\alpha=1}^K)$ solves the same labeled ODE system with permuted initial data. If two labeled configurations have the same initial empirical measure, then because both measures are sums of $K$ equally weighted Dirac masses, one initial configuration is obtained from the other by such a permutation. Therefore they induce the same shared trajectory and the same empirical-measure trajectory.
\end{proof}

\begin{remark}[The closure is not specific to the two-layer proxy]
\label{rem:abstract-closure}
The argument above uses only the factorization
\[
  f_\alpha(x)=f(x;q,z_\alpha),
\]
namely that the output of head $\alpha$ depends on the shared state $q$ and on its own head state $z_\alpha$, but not directly on the states of the other heads. Therefore the same empirical-measure closure applies to any fixed-width shared-backbone architecture of this form, including literal fixed-width BatchEnsemble/TabM variants in which $q$ collects all shared weights and $z_\alpha$ collects all member-specific modulation vectors and readout parameters. Parameter sharing creates a common environment; it does not obstruct the head empirical measure from being the correct closed macroscopic variable.
\end{remark}

\section{Head-Learning Scaling}
\label{sec:scalings}
\subsection{A Unified Head-Learning Scaling}

We write the finite-$K$ system directly in terms of the effective head-learning scale $\lambda_K$. The head-learning scale may depend on $K$ through any sequence
\[
  \lambda_K \to \lambda \in [0,\infty).
\]
To relate $\lambda_K$ to actual optimization, suppose the head-averaged risk \eqref{eq:riskK} is trained by gradient flow with shared-side learning rate $\eta_{q,K}$ and head-side learning rate $\eta_{z,K}$. After rescaling time by $\eta_{q,K}$, the dynamics takes the form
\begin{align}
  \dot q_t^K
  &=
  \frac1K\sum_{\alpha=1}^K \mathcal F(q_t^K,z_{\alpha,t}^K),
  \label{eq:raw-scaled-q}
  \\
  \dot z_{\alpha,t}^K
  &=
  \frac{\eta_{z,K}}{K\eta_{q,K}}\,
  \mathcal G(q_t^K,z_{\alpha,t}^K).
  \label{eq:raw-scaled-z}
\end{align}
Thus $\lambda_K=\eta_{z,K}/(K\eta_{q,K})$. In particular, equal learning rates give $\lambda_K=1/K\to 0$, so the large-$K$ limit freezes the head law at leading order. The present note keeps $\lambda_K\to\lambda$ explicit in order to treat both this degenerate regime ($\lambda=0$) and the nontrivial transport regime ($\lambda>0$), the latter requiring $\eta_{z,K}\sim K\eta_{q,K}$.

The corresponding empirical system is exactly
\begin{align}
  \dot q_t^K
  &=
  \int_{\mathcal Z} \mathcal F(q_t^K,z)\,\nu_t^K(\dd z),
  \label{eq:scaled-q}
  \\
  \partial_t \nu_t^K
  +
  \lambda_K\nabla_z\cdot\bigl(\nu_t^K\,\mathcal G(q_t^K,z)\bigr)
  &=
  0,
  \label{eq:scaled-nu}
\end{align}
which suggests the limit system
\begin{align}
  \dot q_t
  &=
  \int_{\mathcal Z} \mathcal F(q_t,z)\,\nu_t(\dd z),
  \label{eq:limit-q}
  \\
  \partial_t \nu_t
  +
  \lambda\nabla_z\cdot\bigl(\nu_t\,\mathcal G(q_t,z)\bigr)
  &=
  0
  \label{eq:limit-nu}
\end{align}
with initial condition $(q_0,\nu_0)\in\mathcal Q\times\PP_1(\mathcal Z)$.

\begin{theorem}[General head mean-field limit]
\label{thm:head-mf}
Assume that $\mathcal F$ and $\mathcal G$ are globally bounded and globally Lipschitz, and let $\lambda_K\to\lambda\in[0,\infty)$. Then:
\begin{enumerate}[leftmargin=2em]
  \item For every initial condition $(q_0,\nu_0)\in\mathcal Q\times\PP_1(\mathcal Z)$, the coupled ODE-transport system \eqref{eq:limit-q}--\eqref{eq:limit-nu} admits a unique global characteristic solution on every finite time interval.
  \item For every horizon $T<\infty$, there exists $C_T<\infty$ such that if $(q_t,\nu_t)$ and $(\bar q_t,\bar\nu_t)$ solve \eqref{eq:limit-q}--\eqref{eq:limit-nu} with the same $\lambda$, then
  \begin{equation}
    \sup_{t\in[0,T]}
    \bigl(
      \|q_t-\bar q_t\|+\W(\nu_t,\bar\nu_t)
    \bigr)
    \le
    C_T
    \bigl(
      \|q_0-\bar q_0\|+\W(\nu_0,\bar\nu_0)
    \bigr).
    \label{eq:head-mf-stability}
  \end{equation}
  \item Let $(q_t^K,\nu_t^K)$ solve \eqref{eq:scaled-q}--\eqref{eq:scaled-nu}, and assume
  \[
    q_0^K\to q_0,
    \qquad
    \W(\nu_0^K,\nu_0)\to 0.
  \]
  Then for every finite horizon $T$,
  \begin{equation}
    \sup_{t\in[0,T]}
    \bigl(
      \|q_t^K-q_t\|+\W(\nu_t^K,\nu_t)
    \bigr)
    \le
    C_T
    \bigl(
      \|q_0^K-q_0\|+\W(\nu_0^K,\nu_0)+|\lambda_K-\lambda|
    \bigr).
    \label{eq:head-mf-conv}
  \end{equation}
\end{enumerate}
\end{theorem}

\begin{proof}
Let $B_F,B_G<\infty$ be global bounds for $\mathcal F,\mathcal G$, and let $L_F,L_G<\infty$ be global Lipschitz constants.

For a continuous path $q\in C([0,T];\mathcal Q)$ and a scalar $\lambda\ge 0$, let $\Phi_t^{q,\lambda}$ denote the characteristic flow of
\[
  \dot Z_t=\lambda\,\mathcal G(q_t,Z_t).
\]
Because $\mathcal G$ is globally bounded and globally Lipschitz, this flow is globally well defined on every finite interval. Define the pushed-forward measure
\[
  \nu_t^{q,\lambda}:=(\Phi_t^{q,\lambda})_\#\nu_0.
\]
For fixed $\lambda$ and two paths $q,\bar q$, \cref{prop:flow-stability} with identical initial points yields
\begin{equation}
  \W(\nu_t^{q,\lambda},\nu_t^{\bar q,\lambda})
  \le
  e^{\lambda L_G t}\lambda L_G\int_0^t \|q_s-\bar q_s\|\,\dd s.
  \label{eq:abstract-flow-q-diff}
\end{equation}

Define the map $\Gamma$ on $C([0,T];\mathcal Q)$ by
\[
  (\Gamma q)_t
  :=
  q_0 + \int_0^t \int_{\mathcal Z}\mathcal F(q_s,z)\,\nu_s^{q,\lambda}(\dd z)\,\dd s.
\]
For $q,\bar q\in C([0,\tau];\mathcal Q)$ one has
\begin{align*}
  \|(\Gamma q)_t-(\Gamma\bar q)_t\|
  &\le
  \int_0^t
  \Bigl\|
    \int \mathcal F(q_s,z)\,\nu_s^{q,\lambda}(\dd z)
    -
    \int \mathcal F(\bar q_s,z)\,\nu_s^{\bar q,\lambda}(\dd z)
  \Bigr\|
  \dd s
  \\
  &\le
  L_F\int_0^t
  \bigl(
    \|q_s-\bar q_s\|
    +
    \W(\nu_s^{q,\lambda},\nu_s^{\bar q,\lambda})
  \bigr)\dd s
  \\
  &\le
  L_F\tau
  \Bigl(
    1+\lambda L_G \tau e^{\lambda L_G\tau}
  \Bigr)
  \|q-\bar q\|_{L^\infty([0,\tau])}.
\end{align*}
Hence $\Gamma$ is a contraction on $C([0,\tau];\mathcal Q)$ whenever $\tau>0$ is chosen so that the coefficient on the right is $<1$. Banach's fixed-point theorem therefore yields a unique solution on $[0,\tau]$ to \eqref{eq:limit-q}--\eqref{eq:limit-nu}. Since
\[
  \|q_t\|\le \|q_0\|+B_F t,
\]
the same contraction estimate can be iterated finitely many times to cover any finite horizon $[0,T]$. This proves item 1.

For item 2, let $(q_t,\nu_t)$ and $(\bar q_t,\bar\nu_t)$ be two solutions with the same $\lambda$. Let $\pi_0$ be an optimal coupling for $\W(\nu_0,\bar\nu_0)$ and define
\[
  \pi_t:=
  (\Phi_t^{q,\lambda},\Phi_t^{\bar q,\lambda})_\#\pi_0.
\]
Then $\pi_t$ is a coupling of $\nu_t$ and $\bar\nu_t$, so by \cref{prop:flow-stability},
\begin{equation}
  \W(\nu_t,\bar\nu_t)
  \le
  e^{\lambda L_G t}
  \Bigl(
    \W(\nu_0,\bar\nu_0)
    +
    \lambda L_G\int_0^t \|q_s-\bar q_s\|\,\dd s
  \Bigr).
  \label{eq:abstract-flow-stab}
\end{equation}
Also,
\begin{align*}
  \|q_t-\bar q_t\|
  &\le
  \|q_0-\bar q_0\|
  +
  \int_0^t
  \Bigl\|
    \int \mathcal F(q_s,z)\,\nu_s(\dd z)
    -
    \int \mathcal F(\bar q_s,z)\,\bar\nu_s(\dd z)
  \Bigr\|
  \dd s
  \\
  &\le
  \|q_0-\bar q_0\|
  +
  L_F\int_0^t
  \bigl(
    \|q_s-\bar q_s\|+\W(\nu_s,\bar\nu_s)
  \bigr)\dd s.
\end{align*}
Set
\[
  E_t:=\sup_{u\in[0,t]}\bigl(\|q_u-\bar q_u\|+\W(\nu_u,\bar\nu_u)\bigr).
\]
Using \eqref{eq:abstract-flow-stab} in the last bound and then taking suprema gives
\[
  E_t
  \le
  C_T
  \bigl(
    \|q_0-\bar q_0\|+\W(\nu_0,\bar\nu_0)
  \bigr)
  +
  C_T\int_0^t E_s\,\dd s
\]
for a finite constant $C_T$ depending only on $T,\lambda,L_F,L_G$. Gronwall proves \eqref{eq:head-mf-stability}.

For item 3, because $\lambda_K\to\lambda$, the sequence $(\lambda_K)$ is bounded; let $\Lambda:=\sup_K\lambda_K\vee\lambda<\infty$. Let $\pi_0^K$ be an optimal coupling between $\nu_0^K$ and $\nu_0$, and define
\[
  \pi_t^K:=
  (\Phi_t^{q^K,\lambda_K},\Phi_t^{q,\lambda})_\#\pi_0^K.
\]
Then $\pi_t^K$ couples $\nu_t^K$ and $\nu_t$, so \cref{prop:flow-stability} gives
\begin{align}
  \W(\nu_t^K,\nu_t)
  &\le
  e^{\Lambda L_G t}
  \Bigl(
    \W(\nu_0^K,\nu_0)
    +
    \Lambda L_G\int_0^t \|q_s^K-q_s\|\,\dd s
    +
    B_G t\,|\lambda_K-\lambda|
  \Bigr).
  \label{eq:abstract-flow-conv}
\end{align}
Likewise,
\begin{align*}
  \|q_t^K-q_t\|
  &\le
  \|q_0^K-q_0\|
  +
  L_F\int_0^t
  \bigl(
    \|q_s^K-q_s\|+\W(\nu_s^K,\nu_s)
  \bigr)\dd s.
\end{align*}
With
\[
  E_t^K:=\sup_{u\in[0,t]}\bigl(\|q_u^K-q_u\|+\W(\nu_u^K,\nu_u)\bigr),
\]
\eqref{eq:abstract-flow-conv} yields
\[
  E_t^K
  \le
  C_T
  \bigl(
    \|q_0^K-q_0\|+\W(\nu_0^K,\nu_0)+|\lambda_K-\lambda|
  \bigr)
  +
  C_T\int_0^t E_s^K\,\dd s.
\]
Gronwall gives \eqref{eq:head-mf-conv}.
\end{proof}

\begin{proposition}[Characteristic-flow stability with a common environment]
\label{prop:flow-stability}
Assume that $\mathcal G$ is globally bounded by $B_G$ and globally Lipschitz with constant $L_G$. Fix a finite horizon $T<\infty$, a bound $\Lambda<\infty$, two continuous shared-state paths $q,\bar q\in C([0,T];\mathcal Q)$, two scalars $\lambda,\bar\lambda\in[0,\Lambda]$, and initial points $z,\bar z\in\mathcal Z$. Let $\Phi_t^{q,\lambda}(z)$ and $\Phi_t^{\bar q,\bar\lambda}(\bar z)$ denote the corresponding characteristic trajectories:
\[
  \dot Z_t=\lambda\,\mathcal G(q_t,Z_t),
  \qquad
  \dot{\bar Z}_t=\bar\lambda\,\mathcal G(\bar q_t,\bar Z_t).
\]
Then for every $t\in[0,T]$,
\begin{align}
  \|\Phi_t^{q,\lambda}(z)-\Phi_t^{\bar q,\bar\lambda}(\bar z)\|
  \le
  e^{\Lambda L_G t}
  \Bigl(
    \|z-\bar z\|
    +
    \Lambda L_G\int_0^t \|q_s-\bar q_s\|\,\dd s
    +
    B_G t\,|\lambda-\bar\lambda|
  \Bigr).
  \label{eq:flow-stability}
\end{align}
\end{proposition}

\begin{proof}
Write $\Delta_t:=\|\Phi_t^{q,\lambda}(z)-\Phi_t^{\bar q,\bar\lambda}(\bar z)\|$. Then
\[
  \dot \Delta_t
  \le
  \lambda
  \bigl\|
    \mathcal G(q_t,\Phi_t^{q,\lambda}(z))
    -
    \mathcal G(\bar q_t,\Phi_t^{\bar q,\bar\lambda}(\bar z))
  \bigr\|
  +
  |\lambda-\bar\lambda|\,
  \|\mathcal G(\bar q_t,\Phi_t^{\bar q,\bar\lambda}(\bar z))\|.
\]
Using the Lipschitz bound of $\mathcal G$ in $(q,z)$ and the uniform bound $B_G$ gives
\[
  \dot \Delta_t
  \le
  \Lambda L_G \Delta_t
  +
  \Lambda L_G \|q_t-\bar q_t\|
  +
  B_G |\lambda-\bar\lambda|.
\]
Gronwall yields \eqref{eq:flow-stability}.
\end{proof}

\begin{remark}[What kind of mean field this is]
\label{rem:semidirect}
The limit in \cref{thm:head-mf} is not the same object as the width mean-field limit studied in the main draft. It is a \emph{semidirect} head mean field: the head population is described by a transport equation, while all coupling between heads is mediated through the finite-dimensional common environment $q_t$. The heads do not need to move identically, and they do not form an independent-particle system in the classical sense. What makes the limit work is the exact empirical closure from \cref{prop:closure}, not any assumption that the shared particles themselves are exchangeable.
\end{remark}

\begin{corollary}[IID initialization and propagation of chaos]
\label{cor:poc}
Assume the hypotheses of \cref{thm:head-mf}. Suppose that $q_0^K\to q_0$ deterministically and that there exists an iid sequence $(Z_{\alpha,0})_{\alpha\ge 1}$ with common law $\nu_0$ such that, for every $K$,
\[
  z_{\alpha,0}^K=Z_{\alpha,0},
  \qquad
  \alpha=1,\dots,K.
\]
Then almost surely
\[
  \sup_{t\in[0,T]}
  \bigl(
    \|q_t^K-q_t\|+\W(\nu_t^K,\nu_t)
  \bigr)\to 0
  \qquad
  \text{for every }T<\infty.
\]
Moreover, every fixed number of tagged heads becomes asymptotically iid in law, with common marginal given by the transported law $\nu_t$.
\end{corollary}

\begin{proof}
By the strong law for empirical measures in $\W$ on $\PP_1(\mathcal Z)$, the empirical measures
\[
  \nu_0^K=\frac1K\sum_{\alpha=1}^K \delta_{Z_{\alpha,0}}
\]
satisfy
\[
  \W(\nu_0^K,\nu_0)\to 0
  \qquad
  \text{almost surely.}
\]
The first claim therefore follows directly from \cref{thm:head-mf}.

Fix $n\ge 1$ and consider the first $n$ tagged heads. On the same probability space, define
\[
  \bar Z_{i,t}:=\Phi_t^{q,\lambda}(Z_{i,0}),
  \qquad
  i=1,\dots,n,
\]
where $(q_t,\nu_t)$ is the limiting solution. Since the initial variables $(Z_{i,0})_{i=1}^n$ are iid with law $\nu_0$, the limit variables $(\bar Z_{i,t})_{i=1}^n$ are iid with common law $\nu_t$.

For each fixed $i$, \cref{prop:flow-stability} with identical initial point $Z_{i,0}$ gives
\[
  \sup_{t\in[0,T]}
  \|z_{i,t}^K-\bar Z_{i,t}\|
  \le
  e^{\Lambda L_G T}
  \Bigl(
    \Lambda L_G \int_0^T \|q_s^K-q_s\|\,\dd s
    +
    B_G T|\lambda_K-\lambda|
  \Bigr),
\]
where $\Lambda:=\sup_K\lambda_K\vee\lambda<\infty$. The right-hand side converges almost surely to zero by \cref{thm:head-mf}. Hence
\[
  (z_{1,t}^K,\dots,z_{n,t}^K)
  \to
  (\bar Z_{1,t},\dots,\bar Z_{n,t})
  \qquad
  \text{almost surely in }(\mathcal Z)^n
\]
uniformly on $[0,T]$. Therefore the tagged heads converge jointly in law to iid copies with common marginal $\nu_t$.
\end{proof}

\section{Structural Consequences}
\label{sec:consequences}

\subsection{Ensemble Average as a Law-of-Large-Numbers Object}

The most immediate observable in the large-$K$ regime is the ensemble-averaged predictor
\[
  \bar f_t^K(x) := \frac1K\sum_{\alpha=1}^K f(x;q_t^K,z_{\alpha,t}^K).
\]

\begin{corollary}[Deterministic large-ensemble predictor]
\label{cor:predictor}
Under the assumptions of \cref{thm:head-mf}, fix a bounded input set $\mathcal X_\ast\subset\mathcal X$ and a finite horizon $T<\infty$. Assume in addition that $(q,z)\mapsto f(x;q,z)$ is jointly globally Lipschitz on $\mathcal Q\times\mathcal Z$, uniformly over $x\in\mathcal X_\ast$. Then
\[
  \sup_{t\in[0,T]}\sup_{x\in\mathcal X_\ast}
  \left|
    \bar f_t^K(x) - \bar f_t(x)
  \right|
  \to 0,
\]
where
\[
  \bar f_t(x)
  :=
  \int_{\mathcal Z} f(x;q_t,z)\,\nu_t(\dd z).
\]
\end{corollary}

\begin{proof}
The additional hypothesis makes $(q,\nu)\mapsto \int f(x;q,z)\nu(\dd z)$ Lipschitz uniformly over $x\in\mathcal X_\ast$ with respect to $\|q-\bar q\|+\W(\nu,\bar\nu)$. The claim therefore follows directly from \eqref{eq:head-mf-conv}.
\end{proof}

From a practical point of view, however, the averaged predictor is still not the full story. The more informative object is the \emph{predictor law} induced by the head law under the map $z\mapsto f(\cdot;q_t,z)$. We record the state-level no-collapse fact first because it isolates what the transport equation can and cannot do before turning to function-level collapse itself.

\subsection{Bi-Lipschitz Transport and No Finite-Time State Collapse}

In the large-$K$ formulation, the right analogue of headwise collapse is no longer equality of finitely many functions. Instead, it is concentration of the head law onto a single state. Under the bounded-Lipschitz dynamics above, even that state-level collapse cannot occur in finite time unless it is already present initially.

\begin{proposition}[Bi-Lipschitz transport and no finite-time state collapse]
\label{prop:mono}
Assume the hypotheses of \cref{thm:head-mf}. For every finite horizon $T<\infty$ and every $t\in[0,T]$, there exists a bi-Lipschitz map
\[
  \Phi_t:\mathcal Z\to\mathcal Z
\]
such that
\[
  \nu_t = \Phi_{t\#}\nu_0.
\]
Consequently, $\nu_t$ is a Dirac mass if and only if $\nu_0$ is a Dirac mass. In particular, if $\nu_0=\delta_{z_0}$, then $\nu_t=\delta_{z_t}$ where $(q_t,z_t)$ solves
\begin{align}
  \dot q_t &= \mathcal F(q_t,z_t),
  \label{eq:mono-q}
  \\
  \dot z_t &= \lambda\,\mathcal G(q_t,z_t).
  \label{eq:mono-z}
\end{align}
\end{proposition}

\begin{proof}
For the already constructed solution $(q_t,\nu_t)$, let $\Phi_t$ be the characteristic flow of
\[
  \dot Z_t = \lambda\,\mathcal G(q_t,Z_t),
\]
so that $\nu_t=\Phi_{t\#}\nu_0$. Because the vector field is globally Lipschitz in $z$, the characteristic ODE has unique forward and backward solutions on every finite interval. Hence $\Phi_t$ is bijective, and standard Gronwall estimates give bi-Lipschitz bounds for both $\Phi_t$ and $\Phi_t^{-1}$. A pushforward by a bijection is a Dirac mass exactly when the original measure is a Dirac mass. The monomorphic ODE system is just the specialization to $\nu_0=\delta_{z_0}$.
\end{proof}

\paragraph{Interpretation.}
\Cref{prop:mono} is stronger than the earlier monomorphic-stays-monomorphic statement. It shows that, at the level of head states, finite-time exact collapse from dispersed initial data is impossible in the head mean-field regime. If one nevertheless observes function-level collapse of the predictors, that must come from degeneracy of the map
\[
  z \mapsto f(\cdot;q_t,z)
\]
rather than from the transport equation crushing the head law to a point.

\subsection{Predictor-Law Regimes: Disagreement and Bias}

State-level polymorphism is not the same thing as predictor-level diversity. The
practically relevant function-space object is the predictor law obtained by pushing the head law
forward through the map $z\mapsto f(\cdot;q_t,z)$.

\begin{proposition}[Predictor-law collapse criterion and disagreement variance]
\label{prop:disagreement}
Assume the hypotheses of \cref{thm:head-mf}, fix $t\ge 0$, and suppose that
\[
  H_t:\mathcal Z\to L^2(\mathsf P_X),
  \qquad
  H_t(z):=f(\cdot;q_t,z),
\]
belongs to $L^2(\nu_t;L^2(\mathsf P_X))$. Define the predictor mean
\[
  \bar H_t := \int_{\mathcal Z} H_t(z)\,\nu_t(\dd z)
\]
and the disagreement energy
\[
  \mathcal D_t
  :=
  \iint_{\mathcal Z\times\mathcal Z}
  \|H_t(z)-H_t(z')\|_{L^2(\mathsf P_X)}^2\,
  \nu_t(\dd z)\nu_t(\dd z').
\]
Define also the predictor law
\[
  \Pi_t := (H_t)_\# \nu_t \in \PP\bigl(L^2(\mathsf P_X)\bigr).
\]
Then
\begin{equation}
  \mathcal D_t
  =
  2\int_{\mathcal Z}
  \|H_t(z)-\bar H_t\|_{L^2(\mathsf P_X)}^2\,
  \nu_t(\dd z).
  \label{eq:disagreement-variance}
\end{equation}
Moreover, the following are equivalent:
\begin{enumerate}[leftmargin=2em]
  \item $\Pi_t$ is a Dirac mass;
  \item $\mathcal D_t=0$;
  \item $H_t(z)=\bar H_t$ for $\nu_t$-almost every $z$.
\end{enumerate}
\end{proposition}

\begin{proof}
Let $Z,Z'$ be iid random variables with law $\nu_t$. Since $H_t(Z)$ is a
square-integrable random element of the Hilbert space $L^2(\mathsf P_X)$,
\[
  \mathcal D_t
  =
  \E\|H_t(Z)-H_t(Z')\|_{L^2(\mathsf P_X)}^2.
\]
Expanding the square in a Hilbert space and using the independence and identical
distribution of $Z$ and $Z'$ gives
\begin{align*}
  \mathcal D_t
  &=
  \E\|H_t(Z)\|_{L^2(\mathsf P_X)}^2
  +
  \E\|H_t(Z')\|_{L^2(\mathsf P_X)}^2
  -
  2\bigl\langle \E H_t(Z), \E H_t(Z') \bigr\rangle_{L^2(\mathsf P_X)}
  \\
  &=
  2\E\|H_t(Z)-\bar H_t\|_{L^2(\mathsf P_X)}^2,
\end{align*}
which is \eqref{eq:disagreement-variance}. The characterization of the case
$\mathcal D_t=0$ is the standard fact that a nonnegative square-integrable
random variable has zero expectation if and only if it vanishes almost surely.
This is equivalent to $H_t(Z)$ being almost surely constant, which in turn is
equivalent to the pushforward law $\Pi_t$ being a Dirac mass.
\end{proof}

\begin{corollary}[Injectivity prevents function-level collapse]
\label{cor:injective}
Assume the hypotheses of \cref{prop:disagreement}, and suppose that there exists
$c_t>0$ such that
\[
  \|H_t(z)-H_t(z')\|_{L^2(\mathsf P_X)}
  \ge
  c_t\|z-z'\|
  \qquad
  \forall z,z'\in\operatorname{supp}\nu_t.
\]
Let
\[
  \bar z_t := \int_{\mathcal Z} z\,\nu_t(\dd z),
  \qquad
  V_t^2 := \int_{\mathcal Z}\|z-\bar z_t\|^2\,\nu_t(\dd z).
\]
Then
\begin{equation}
  \mathcal D_t \ge 2c_t^2 V_t^2.
  \label{eq:disagreement-lower}
\end{equation}
In particular, $\Pi_t$ is a Dirac mass if and only if $\nu_t$ is a Dirac mass.
\end{corollary}

\begin{proof}
Let $Z,Z'$ be iid with law $\nu_t$. The lower Lipschitz bound gives
\[
  \mathcal D_t
  =
  \E\|H_t(Z)-H_t(Z')\|_{L^2(\mathsf P_X)}^2
  \ge
  c_t^2 \E\|Z-Z'\|^2
  =
  2c_t^2 V_t^2.
\]
If $\Pi_t$ is a Dirac mass, then \eqref{eq:disagreement-lower} implies
$V_t=0$, hence $\nu_t$ is a Dirac mass. The converse is immediate.
\end{proof}

\begin{corollary}[Lipschitz sensitivity controls disagreement and Jensen bias]
\label{cor:bias}
Assume the hypotheses of \cref{prop:disagreement}, and in addition assume that
\[
  \int_{\mathcal Z}\|z\|^2\,\nu_t(\dd z)<\infty.
\]
Let
\[
  \bar z_t := \int_{\mathcal Z} z\,\nu_t(\dd z),
  \qquad
  V_t^2 := \int_{\mathcal Z}\|z-\bar z_t\|^2\,\nu_t(\dd z).
\]
If $H_t$ is $L_t$-Lipschitz on $\operatorname{supp}\nu_t$, then
\begin{align}
  \mathcal D_t &\le 2L_t^2 V_t^2,
  \label{eq:disagreement-lip}
  \\
  \|\bar H_t-H_t(\bar z_t)\|_{L^2(\mathsf P_X)}
  &\le
  L_t V_t.
  \label{eq:bias-lip}
\end{align}
If moreover $H_t$ is affine on the convex hull of $\operatorname{supp}\nu_t$,
then
\begin{equation}
  \bar H_t = H_t(\bar z_t).
  \label{eq:bias-affine}
\end{equation}
\end{corollary}

\begin{proof}
Let $Z,Z'$ be iid with law $\nu_t$. By the Lipschitz property of $H_t$,
\[
  \|H_t(Z)-H_t(Z')\|_{L^2(\mathsf P_X)}
  \le
  L_t\|Z-Z'\|.
\]
Squaring and taking expectations yields
\[
  \mathcal D_t
  \le
  L_t^2 \E\|Z-Z'\|^2.
\]
Since $Z,Z'$ are iid,
\[
  \E\|Z-Z'\|^2
  =
  2\E\|Z-\bar z_t\|^2
  =
  2V_t^2,
\]
which proves \eqref{eq:disagreement-lip}. Likewise,
\[
  \|\bar H_t-H_t(\bar z_t)\|_{L^2(\mathsf P_X)}
  =
  \left\|
    \int_{\mathcal Z}\bigl(H_t(z)-H_t(\bar z_t)\bigr)\,\nu_t(\dd z)
  \right\|_{L^2(\mathsf P_X)}
  \le
  \int_{\mathcal Z}\|H_t(z)-H_t(\bar z_t)\|_{L^2(\mathsf P_X)}\,\nu_t(\dd z)
\]
by Minkowski. The Lipschitz bound gives
\[
  \|\bar H_t-H_t(\bar z_t)\|_{L^2(\mathsf P_X)}
  \le
  L_t \int_{\mathcal Z}\|z-\bar z_t\|\,\nu_t(\dd z)
  \le
  L_t V_t,
\]
proving \eqref{eq:bias-lip}. Finally, if $H_t$ is affine on
$\operatorname{co}(\operatorname{supp}\nu_t)$, then $\bar z_t$ belongs to that
convex hull and linearity of Bochner integration gives \eqref{eq:bias-affine}.
\end{proof}

\paragraph{Regime taxonomy.}
\Cref{prop:mono,prop:disagreement,cor:injective,cor:bias} suggest three distinct levels of
regime classification.
\begin{enumerate}[leftmargin=2em]
  \item \emph{Function regime}: the predictor law $\Pi_t$ is monomorphic if it is a Dirac mass and polymorphic otherwise; quantitatively, the ensemble is functionally collective when $\mathcal D_t$ is small and functionally diverse when $\mathcal D_t$ is appreciable. By \cref{prop:disagreement}, this is determined by the predictor law rather than by $\nu_t$ alone.
  \item \emph{State regime}: $\nu_t$ is monomorphic if it is a Dirac mass and polymorphic otherwise. By \cref{prop:mono}, finite-time state collapse cannot occur from dispersed initial data, but by \cref{prop:disagreement} this need not decide the function regime unless the map $H_t$ is nondegenerate on $\operatorname{supp}\nu_t$.
  \item \emph{Bias regime}: the ensemble is bias-negligible when $\|\bar H_t-H_t(\bar z_t)\|$ is small and bias-relevant when this Jensen gap is non-negligible. By \cref{cor:bias}, small Lipschitz sensitivity or approximate affinity of $H_t$ suppresses this bias.
\end{enumerate}
Thus a head population may be state-polymorphic but still functionally
collective if the map $z\mapsto H_t(z)$ is nearly constant on the support of
$\nu_t$. Conversely, if $H_t$ is injective or quantitatively nondegenerate on
that support, then state polymorphism forces function polymorphism by
\cref{cor:injective}. A functionally diverse ensemble can still be
bias-negligible if $H_t$ varies along the support of $\nu_t$ but remains
approximately affine.

\subsection{Finite-Horizon Moment Bounds for the Concrete Proxy}

The main gap identified earlier was that the literal two-layer proxy does not satisfy the global bounded-Lipschitz assumptions of \cref{thm:head-mf}. The next proposition shows that this obstruction is milder than it first appears: despite the lack of global boundedness of the raw drift, the concrete finite-$K$ dynamics obey a closed finite-horizon estimate in terms of second moments of the head population.

\begin{proposition}[Uniform finite-horizon moment bounds for the two-layer proxy]
\label{prop:concrete-apriori}
Assume that the input support is bounded,
\[
  \|x\|\le B_X
  \qquad
  \mathsf P\text{-almost surely},
\]
that $\sigma,\sigma'$ are bounded, and that
\[
  |\partial_1\ell(u,y)|\le L_\ell
  \qquad
  \forall (u,y)\in\R\times\mathcal Y.
\]
Let $\Lambda_h<\infty$ and suppose $0\le \lambda_K\le \Lambda_h$ for all $K$. For the finite-$K$ system \eqref{eq:general-q}--\eqref{eq:general-z}, write
\[
  q_t^K=(w_t^K,b_t^K)\in\R^{md}\times\R^m,
  \qquad
  z_{\alpha,t}^K=(a_{\alpha,t}^K,r_{\alpha,t}^K)\in\R^m\times\R^d,
\]
and define the aggregate quantities
\begin{align}
  A_t^K
  &:=
  \biggl(
    \frac1K\sum_{\alpha=1}^K \|a_{\alpha,t}^K\|^2
  \biggr)^{1/2},
  \label{eq:A2-def}
  \\
  R_t^K
  &:=
  \biggl(
    \frac1K\sum_{\alpha=1}^K \|r_{\alpha,t}^K\|^2
  \biggr)^{1/2},
  \label{eq:R2-def}
  \\
  W_t^K &:= \|w_t^K\|,
  \qquad
  B_t^K := \|b_t^K\|.
  \label{eq:WB-def}
\end{align}
Then there exist constants
\[
  c_a := \frac{L_\ell\|\sigma\|_\infty}{\sqrt m},
  \qquad
  c_b := \frac{L_\ell\|\sigma'\|_\infty}{m},
  \qquad
  c_\ast := \frac{L_\ell\|\sigma'\|_\infty B_X}{m}
\]
such that for every $t\ge 0$,
\begin{align}
  A_t^K
  &\le
  A_0^K + \Lambda_h c_a t,
  \label{eq:concrete-A-bound}
  \\
  B_t^K
  &\le
  B_0^K + c_b \int_0^t A_s^K\,\dd s,
  \label{eq:concrete-B-bound}
  \\
  W_t^K
  &\le
  W_0^K + c_\ast \int_0^t A_s^K R_s^K\,\dd s,
  \label{eq:concrete-W-bound}
  \\
  R_t^K
  &\le
  R_0^K + \Lambda_h c_\ast \int_0^t A_s^K W_s^K\,\dd s.
  \label{eq:concrete-R-bound}
\end{align}
Consequently, with $U_t^K:=W_t^K+R_t^K$,
\begin{equation}
  U_t^K
  \le
  U_0^K
  \exp\!\biggl(
    c_\ast(1+\Lambda_h)
    \int_0^t A_s^K\,\dd s
  \biggr),
  \label{eq:concrete-U-bound}
\end{equation}
and therefore
\begin{equation}
  U_t^K
  \le
  U_0^K
  \exp\!\biggl(
    c_\ast(1+\Lambda_h)
    \Bigl(
      A_0^K t + \frac12 \Lambda_h c_a t^2
    \Bigr)
  \biggr).
  \label{eq:concrete-U-explicit}
\end{equation}
In particular, if
\[
  \sup_K\bigl(A_0^K+R_0^K+W_0^K+B_0^K\bigr)<\infty,
\]
then for every finite horizon $T<\infty$,
\[
  \sup_K\sup_{t\in[0,T]}
  \bigl(
    A_t^K+R_t^K+W_t^K+B_t^K
  \bigr)
  <\infty.
\]
\end{proposition}

\begin{proof}
For a single head state $z=(a,r)$ and shared state $q=(w,b)$, the gradients of \eqref{eq:generic-output} satisfy
\begin{align}
  \|\nabla_a f(x;q,z)\|
  &\le
  \frac{\|\sigma\|_\infty}{\sqrt m},
  \label{eq:grad-a-bound}
  \\
  \|\nabla_b f(x;q,z)\|
  &\le
  \frac{\|\sigma'\|_\infty}{m}\|a\|,
  \label{eq:grad-b-bound}
  \\
  \|\nabla_w f(x;q,z)\|
  &\le
  \frac{\|\sigma'\|_\infty B_X}{m}\|a\|\,\|r\|,
  \label{eq:grad-w-bound}
  \\
  \|\nabla_r f(x;q,z)\|
  &\le
  \frac{\|\sigma'\|_\infty B_X}{m}\|a\|\,\|w\|.
  \label{eq:grad-r-bound}
\end{align}
The first two are immediate. For \eqref{eq:grad-w-bound}, the $j$th block equals
\[
  \frac{a_j}{m}\sigma'(w_j^\top(r\odot x)+b_j)(r\odot x),
\]
so
\[
  \|\nabla_w f\|^2
  \le
  \frac{\|\sigma'\|_\infty^2}{m^2}
  \sum_{j=1}^m a_j^2 \|r\odot x\|^2
  \le
  \frac{\|\sigma'\|_\infty^2 B_X^2}{m^2}
  \|a\|^2\|r\|^2.
\]
The bound \eqref{eq:grad-r-bound} follows similarly from
\[
  \nabla_r f(x;q,z)
  =
  \frac1m\sum_{j=1}^m
  a_j \sigma'(w_j^\top(r\odot x)+b_j)(w_j\odot x)
\]
and Cauchy--Schwarz.

Multiplying the gradient bounds by $L_\ell$ yields the same estimates for the drift magnitudes. The $a$-bound gives
\[
  \|\dot a_{\alpha,t}^K\|
  \le
  \lambda_K c_a
  \le
  \Lambda_h c_a.
\]
Applying Minkowski to the vector
\[
  \bigl(
    \|a_{1,t}^K\|,\dots,\|a_{K,t}^K\|
  \bigr)\in\R^K
\]
gives \eqref{eq:concrete-A-bound}. Next,
\[
  \| \dot b_t^K \|
  \le
  \frac{L_\ell\|\sigma'\|_\infty}{m}
  \frac1K\sum_{\alpha=1}^K \|a_{\alpha,t}^K\|
  \le
  c_b A_t^K,
\]
which yields \eqref{eq:concrete-B-bound}. Likewise,
\[
  \| \dot w_t^K \|
  \le
  \frac{L_\ell\|\sigma'\|_\infty B_X}{m}
  \frac1K\sum_{\alpha=1}^K
  \|a_{\alpha,t}^K\|\,\|r_{\alpha,t}^K\|
  \le
  c_\ast A_t^K R_t^K
\]
by Cauchy--Schwarz, proving \eqref{eq:concrete-W-bound}.

For the $r$-moment, differentiate the empirical second moment:
\[
  \frac{\dd}{\dd t}\frac1K\sum_{\alpha=1}^K \frac12 \|r_{\alpha,t}^K\|^2
  \le
  \lambda_K c_\ast W_t^K
  \frac1K\sum_{\alpha=1}^K
  \|a_{\alpha,t}^K\|\,\|r_{\alpha,t}^K\|
  \le
  \Lambda_h c_\ast W_t^K A_t^K R_t^K.
\]
Using the standard regularization argument for $\sqrt{\varepsilon+(R_t^K)^2}$ and then sending $\varepsilon\downarrow 0$ yields \eqref{eq:concrete-R-bound}. Summing \eqref{eq:concrete-W-bound} and \eqref{eq:concrete-R-bound} gives
\[
  \frac{\dd}{\dd t}U_t^K
  \le
  c_\ast(1+\Lambda_h)A_t^K U_t^K
\]
in integral form, so Gronwall yields \eqref{eq:concrete-U-bound}. Finally,
\[
  \int_0^t A_s^K\,\dd s
  \le
  A_0^K t + \frac12 \Lambda_h c_a t^2,
\]
which gives \eqref{eq:concrete-U-explicit}. The final uniform finite-horizon bound follows immediately.
\end{proof}


\begin{proposition}[Bounded-region Lipschitz structure of the concrete proxy]
\label{prop:moment-local}
Assume the hypotheses of \cref{prop:concrete-apriori}, and in addition assume that $\sigma'$ is Lipschitz on $\R$ and that $\partial_1\ell(\cdot,y)$ is Lipschitz uniformly in $y\in\mathcal Y$. Define
\[
  K_{R_q,R_z}
  :=
  \bigl\{
    (q,z)\in\mathcal Q\times\mathcal Z:
    \|q\|\le R_q,\ \|z\|\le R_z
  \bigr\}.
\]
Then for every pair of bounds $R_q,R_z<\infty$, there exists a constant $C_{R_q,R_z}<\infty$ such that for all $(q,z),(\bar q,\bar z)\in K_{R_q,R_z}$,
\begin{align}
  \|\mathcal F(q,z)\|+\|\mathcal G(q,z)\|
  &\le
  C_{R_q,R_z},
  \label{eq:concrete-local-bdd}
  \\
  \|\mathcal F(q,z)-\mathcal F(\bar q,\bar z)\|
  +
  \|\mathcal G(q,z)-\mathcal G(\bar q,\bar z)\|
  &\le
  C_{R_q,R_z}
  \bigl(
    \|q-\bar q\|+\|z-\bar z\|
  \bigr).
  \label{eq:concrete-local-lip}
\end{align}
\end{proposition}

\begin{proof}
Write $z=(a,r)$, $\bar z=(\bar a,\bar r)$ and $q=(w,b)$, $\bar q=(\bar w,\bar b)$. For fixed $x$, let
\[
  u_j:=w_j^\top(r\odot x)+b_j,
  \qquad
  \bar u_j:=\bar w_j^\top(\bar r\odot x)+\bar b_j.
\]
Since $\|x\|\le B_X$, $\|q\|,\|\bar q\|\le R_q$, and $\|z\|,\|\bar z\|\le R_z$, one has
\begin{align}
  |u_j-\bar u_j|
  &\le
  \|w_j-\bar w_j\|\,\|r\odot x\|
  +
  \|\bar w_j\|\,\|(r-\bar r)\odot x\|
  +
  |b_j-\bar b_j|
  \notag\\
  &\le
  B_X R_z \|w_j-\bar w_j\|
  +
  B_X R_q \|r-\bar r\|
  +
  |b_j-\bar b_j|.
  \label{eq:u-diff-bound}
\end{align}
Summing over $j$ and using $\|v\|_1\le \sqrt m\,\|v\|$ yields
\begin{equation}
  \|u-\bar u\|_1
  \le
  C_{R_q,R_z}
  \bigl(
    \|q-\bar q\|+\|z-\bar z\|
  \bigr).
  \label{eq:u-l1-bound}
\end{equation}

Because $\sigma$ is bounded with bounded derivative and $\sigma'$ is Lipschitz, the maps $\sigma$ and $\sigma'$ are globally Lipschitz. Using \eqref{eq:u-l1-bound} and the explicit formulas
\[
  f(x;q,z)
  =
  \frac1m\sum_{j=1}^m a_j \sigma(u_j),
  \qquad
  \nabla_a f(x;q,z)
  =
  \frac1m\bigl(\sigma(u_j)\bigr)_{j=1}^m,
\]
\[
  \nabla_b f(x;q,z)
  =
  \frac1m\bigl(a_j\sigma'(u_j)\bigr)_{j=1}^m,
  \qquad
  \nabla_w f(x;q,z)
  =
  \frac1m\bigl(a_j\sigma'(u_j)(r\odot x)\bigr)_{j=1}^m,
\]
\[
  \nabla_r f(x;q,z)
  =
  \frac1m\sum_{j=1}^m a_j \sigma'(u_j)(w_j\odot x),
\]
one obtains a constant $C_{R_q,R_z}$ such that for all $x$ in the support of $\mathsf P_X$,
\begin{align}
  |f(x;q,z)|+
  \|\nabla_q f(x;q,z)\|+
  \|\nabla_z f(x;q,z)\|
  &\le
  C_{R_q,R_z},
  \label{eq:f-grad-local-bdd}
  \\
  |f(x;q,z)-f(x;\bar q,\bar z)|
  +
  \|\nabla_q f(x;q,z)-\nabla_q f(x;\bar q,\bar z)\|
  +
  \|\nabla_z f(x;q,z)-\nabla_z f(x;\bar q,\bar z)\|
  &\le
  C_{R_q,R_z}
  \bigl(
    \|q-\bar q\|+\|z-\bar z\|
  \bigr).
  \label{eq:f-grad-local-lip}
\end{align}
Indeed, every term above is a finite sum of products of bounded coordinates of $a,r,w,b$ with either $\sigma(u_j)$ or $\sigma'(u_j)$, and the only nonlinear dependence enters through $u_j$, whose Lipschitz control is given by \eqref{eq:u-l1-bound}.

Let $L_{\ell,\mathrm{Lip}}$ denote a uniform Lipschitz constant for $\partial_1\ell(\cdot,y)$ in the first argument. Then \eqref{eq:f-grad-local-bdd} gives
\[
  \|\mathcal F(q,z)\|+\|\mathcal G(q,z)\|
  \le
  L_\ell C_{R_q,R_z},
\]
which is \eqref{eq:concrete-local-bdd}. For the Lipschitz estimate, write
\begin{align*}
  \mathcal F(q,z)-\mathcal F(\bar q,\bar z)
  &=
  -\E
  \Bigl[
    \bigl(
      \partial_1\ell(f(x;q,z),y)
      -
      \partial_1\ell(f(x;\bar q,\bar z),y)
    \bigr)
    \nabla_q f(x;q,z)
  \Bigr]
  \\
  &\quad
  -\E
  \Bigl[
    \partial_1\ell(f(x;\bar q,\bar z),y)
    \bigl(
      \nabla_q f(x;q,z)
      -
      \nabla_q f(x;\bar q,\bar z)
    \bigr)
  \Bigr].
\end{align*}
Taking norms, using the boundedness of $\partial_1\ell$, the Lipschitz property of $\partial_1\ell(\cdot,y)$, and \eqref{eq:f-grad-local-bdd}--\eqref{eq:f-grad-local-lip}, we obtain
\[
  \|\mathcal F(q,z)-\mathcal F(\bar q,\bar z)\|
  \le
  C_{R_q,R_z}
  \bigl(
    \|q-\bar q\|+\|z-\bar z\|
  \bigr).
\]
The same argument with $\nabla_z f$ in place of $\nabla_q f$ gives the corresponding bound for $\mathcal G$. Summing the two estimates proves \eqref{eq:concrete-local-lip}.
\end{proof}

\begin{theorem}[Concrete large-$K$ limit for the two-layer proxy on finite horizons]
\label{thm:concrete-head-mf}
Assume the hypotheses of \cref{prop:moment-local}, let $\lambda_K\to\lambda\in[0,\infty)$ with $\sup_K\lambda_K<\infty$, and suppose
\[
  q_0^K\to q_0,
  \qquad
  \W(\nu_0^K,\nu_0)\to 0.
\]
Assume moreover that there exists $R_0<\infty$ such that
\[
  \|q_0\|\le R_0,
  \qquad
  \sup_K\|q_0^K\|\le R_0,
  \qquad
  \operatorname{supp}\nu_0\subset B_{\mathcal Z}(R_0),
  \qquad
  \operatorname{supp}\nu_0^K\subset B_{\mathcal Z}(R_0)
  \quad
  \forall K,
\]
where $B_{\mathcal Z}(R):=\{z\in\mathcal Z:\|z\|\le R\}$. Then for every finite horizon $T<\infty$:
\begin{enumerate}[leftmargin=2em]
  \item the finite-$K$ concrete system is globally well posed on $[0,T]$;
  \item there exists a unique limit pair $(q_t,\nu_t)_{t\in[0,T]}$ with
  \[
    q\in C([0,T];\mathcal Q),
    \qquad
    \nu\in C([0,T];\PP_1(\mathcal Z)),
  \]
  solving
  \begin{align}
    \dot q_t
    &=
    \int_{\mathcal Z}\mathcal F(q_t,z)\,\nu_t(\dd z),
    \label{eq:concrete-limit-q}
    \\
    \partial_t \nu_t
    +
    \lambda\nabla_z\cdot\bigl(\nu_t\,\mathcal G(q_t,z)\bigr)
    &=
    0
    \label{eq:concrete-limit-nu}
  \end{align}
  in the characteristic sense;
  \item one has convergence
  \[
    \sup_{t\in[0,T]}
    \bigl(
      \|q_t^K-q_t\|+\W(\nu_t^K,\nu_t)
    \bigr)
    \to 0.
  \]
\end{enumerate}
\end{theorem}

\begin{proof}
Set
\[
  \Lambda_h:=\sup_K\lambda_K\vee\lambda<\infty.
\]
Because the concrete vector field is locally Lipschitz on the finite-dimensional space $\mathcal Q\times\mathcal Z^K$, the finite-$K$ system is locally well posed. We first show that no trajectory can blow up on $[0,T]$.

For the empirical quantities from \cref{prop:concrete-apriori}, the bounded-support assumption implies
\[
  A_0^K\le R_0,
  \qquad
  R_0^K\le R_0,
  \qquad
  W_0^K+B_0^K\le \sqrt2\,R_0.
\]
Applying \cref{prop:concrete-apriori} therefore yields constants $R_{q,T},R_{z,T}^{\mathrm{mom}}<\infty$, depending only on $T$, $R_0$, $\Lambda_h$, and the model constants, such that
\begin{equation}
  \sup_K\sup_{t\in[0,T]}
  \bigl(
    W_t^K+B_t^K
  \bigr)
  \le
  R_{q,T}.
  \label{eq:concrete-q-ball}
\end{equation}
\begin{equation}
  \sup_K\sup_{t\in[0,T]}
  \bigl(
    A_t^K+R_t^K
  \bigr)
  \le
  R_{z,T}^{\mathrm{mom}}.
  \label{eq:concrete-z-moment-ball}
\end{equation}
In particular,
\[
  \|q_t^K\|
  \le
  W_t^K+B_t^K
  \le
  R_{q,T}
  \qquad
  \forall t\in[0,T],\ \forall K.
\]
Likewise, for each head $\alpha$,
\[
  \|\dot a_{\alpha,t}^K\|
  \le
  \lambda_K c_a
  \le
  \Lambda_h c_a,
\]
so
\begin{equation}
  \|a_{\alpha,t}^K\|
  \le
  R_0+\Lambda_h c_a T
  =:
  R_{a,T}.
  \label{eq:concrete-a-support}
\end{equation}
Using \eqref{eq:grad-r-bound} and \eqref{eq:concrete-q-ball},
\[
  \|\dot r_{\alpha,t}^K\|
  \le
  \lambda_K c_\ast \|a_{\alpha,t}^K\|\,\|w_t^K\|
  \le
  \Lambda_h c_\ast R_{a,T} R_{q,T},
\]
and therefore
\begin{equation}
  \|r_{\alpha,t}^K\|
  \le
  R_0+\Lambda_h c_\ast R_{a,T}R_{q,T}T
  =:
  R_{r,T}.
  \label{eq:concrete-r-support}
\end{equation}
With
\[
  R_{z,T}:=\sqrt{R_{a,T}^2+R_{r,T}^2},
\]
\eqref{eq:concrete-a-support}--\eqref{eq:concrete-r-support} imply
\begin{equation}
  \operatorname{supp}\nu_t^K\subset B_{\mathcal Z}(R_{z,T})
  \qquad
  \forall t\in[0,T],\ \forall K.
  \label{eq:concrete-z-ball}
\end{equation}
Thus the finite-$K$ solution cannot leave a compact region on $[0,T]$, proving item 1.

Next fix the closed Euclidean balls
\[
  B_{\mathcal Q}:=\{q\in\mathcal Q:\|q\|\le R_{q,T}\},
  \qquad
  B_{\mathcal Z}:=B_{\mathcal Z}(R_{z,T}),
\]
and let
\[
  \Pi_{\mathcal Q}:\mathcal Q\to B_{\mathcal Q},
  \qquad
  \Pi_{\mathcal Z}:\mathcal Z\to B_{\mathcal Z}
\]
denote the metric projections. Since these balls are closed and convex, both projections are $1$-Lipschitz. By \cref{prop:moment-local}, there exists a constant $L_T<\infty$ such that $\mathcal F$ and $\mathcal G$ are bounded by $L_T$ and $L_T$-Lipschitz on $B_{\mathcal Q}\times B_{\mathcal Z}$. Define the globally bounded-Lipschitz extensions
\[
  \widetilde{\mathcal F}(q,z)
  :=
  \mathcal F(\Pi_{\mathcal Q}q,\Pi_{\mathcal Z}z),
  \qquad
  \widetilde{\mathcal G}(q,z)
  :=
  \mathcal G(\Pi_{\mathcal Q}q,\Pi_{\mathcal Z}z).
\]
Because the projections are $1$-Lipschitz, $\widetilde{\mathcal F}$ and $\widetilde{\mathcal G}$ are globally bounded and globally Lipschitz on $\mathcal Q\times\mathcal Z$. Moreover, by \eqref{eq:concrete-q-ball} and \eqref{eq:concrete-z-ball}, every finite-$K$ trajectory of the original concrete system stays inside $B_{\mathcal Q}\times B_{\mathcal Z}$, so it also solves the extended system
\begin{align}
  \dot q_t^K
  &=
  \int_{\mathcal Z}\widetilde{\mathcal F}(q_t^K,z)\,\nu_t^K(\dd z),
  \label{eq:concrete-extended-q}
  \\
  \partial_t \nu_t^K
  +
  \lambda_K\nabla_z\cdot\bigl(
    \nu_t^K\widetilde{\mathcal G}(q_t^K,z)
  \bigr)
  &=
  0.
  \label{eq:concrete-extended-nu}
\end{align}

Applying \cref{thm:head-mf} to the extended fields with initial data $(q_0,\nu_0)$ yields a unique global characteristic solution $(q_t,\nu_t)_{t\in[0,T]}$ and convergence
\begin{equation}
  \sup_{t\in[0,T]}
  \bigl(
    \|q_t^K-q_t\|+\W(\nu_t^K,\nu_t)
  \bigr)
  \to 0
  \qquad
  \text{for the extended system.}
  \label{eq:concrete-extended-conv}
\end{equation}
Since each $q_t^K$ belongs to the closed ball $B_{\mathcal Q}$, the uniform convergence in \eqref{eq:concrete-extended-conv} implies
\[
  q_t\in B_{\mathcal Q}
  \qquad
  \forall t\in[0,T].
\]
Likewise, \eqref{eq:concrete-z-ball} shows that each $\nu_t^K$ is supported in the closed ball $B_{\mathcal Z}$; because $\W$ convergence implies weak convergence, the limit measure $\nu_t$ is also supported in $B_{\mathcal Z}$ for every $t$. Hence, along the limit trajectory,
\[
  \widetilde{\mathcal F}(q_t,z)=\mathcal F(q_t,z),
  \qquad
  \widetilde{\mathcal G}(q_t,z)=\mathcal G(q_t,z)
  \qquad
  \text{for } z\in\operatorname{supp}\nu_t.
\]
Therefore $(q_t,\nu_t)$ solves the original concrete system \eqref{eq:concrete-limit-q}--\eqref{eq:concrete-limit-nu}. This proves existence in item 2, and \eqref{eq:concrete-extended-conv} gives item 3.

It remains to prove uniqueness for the original concrete limit system. Let $(\bar q_t,\bar\nu_t)_{t\in[0,T]}$ be another characteristic solution of \eqref{eq:concrete-limit-q}--\eqref{eq:concrete-limit-nu} with the same initial condition. Define
\[
  \bar A_t
  :=
  \biggl(
    \int_{\mathcal Z}\|a\|^2\,\bar\nu_t(\dd z)
  \biggr)^{1/2},
  \qquad
  \bar R_t
  :=
  \biggl(
    \int_{\mathcal Z}\|r\|^2\,\bar\nu_t(\dd z)
  \biggr)^{1/2}.
\]
Replacing the empirical averages in the proof of \cref{prop:concrete-apriori} by integrals against $\bar\nu_t$ shows that the same closed differential inequalities hold for $(\bar A_t,\bar R_t,\|\bar w_t\|,\|\bar b_t\|)$. Hence the same Gronwall argument yields
\[
  \|\bar q_t\|\le R_{q,T}
  \qquad
  \forall t\in[0,T].
\]
By the definition of characteristic solution, for each $t$ there exists a Borel map $\bar Z_t:\mathcal Z\to\mathcal Z$ such that
\[
  \bar\nu_t=(\bar Z_t)_\#\nu_0,
\]
and for $\nu_0$-almost every $z_0=(a_0,r_0)\in\mathcal Z$, the path $\bar Z_t(z_0)=(\bar a_t(z_0),\bar r_t(z_0))$ satisfies
\[
  \partial_t \bar Z_t(z_0)
  =
  \lambda \mathcal G(\bar q_t,\bar Z_t(z_0)),
  \qquad
  \bar Z_0(z_0)=z_0.
\]
Because $\operatorname{supp}\nu_0\subset B_{\mathcal Z}(R_0)$, one has $\|a_0\|\le R_0$ and $\|r_0\|\le R_0$ for $\nu_0$-almost every $z_0=(a_0,r_0)$. Along such a characteristic,
\[
  \|\partial_t \bar a_t(z_0)\|
  \le
  \lambda c_a
  \le
  \Lambda_h c_a,
\]
so
\[
  \|\bar a_t(z_0)\|
  \le
  R_0+\Lambda_h c_a T
  =
  R_{a,T}.
\]
Using \eqref{eq:grad-r-bound} and $\|\bar w_t\|\le \|\bar q_t\|\le R_{q,T}$,
\[
  \|\partial_t \bar r_t(z_0)\|
  \le
  \lambda c_\ast \|\bar a_t(z_0)\|\,\|\bar w_t\|
  \le
  \Lambda_h c_\ast R_{a,T}R_{q,T},
\]
and therefore
\[
  \|\bar r_t(z_0)\|
  \le
  R_0+\Lambda_h c_\ast R_{a,T}R_{q,T}T
  =
  R_{r,T}.
\]
Consequently,
\[
  \bar\nu_t\bigl(B_{\mathcal Z}(R_{z,T})\bigr)=1
  \qquad
  \forall t\in[0,T].
\]
Since $B_{\mathcal Z}(R_{z,T})$ is closed, this implies
\[
  \operatorname{supp}\bar\nu_t\subset B_{\mathcal Z}
  \qquad
  \forall t\in[0,T].
\]
Hence $(\bar q_t,\bar\nu_t)$ also solves the extended system driven by $\widetilde{\mathcal F},\widetilde{\mathcal G}$. By the uniqueness part of \cref{thm:head-mf} for the extended system, $(\bar q_t,\bar\nu_t)=(q_t,\nu_t)$ on $[0,T]$. This completes the proof.
\end{proof}

\paragraph{Interpretation.}
\Cref{thm:concrete-head-mf} is now a rigorous finite-horizon large-$K$ theorem for the literal two-layer BatchEnsemble-like proxy. The price of rigor is that the theorem is stated for uniformly bounded-support initial head laws, which is exactly the regime in which the compact-region localization closes. This is already a practically meaningful regime for bounded shared initializations such as Xavier/He uniform together with bounded head-specific parameters, for instance bounded output-head initializations for $a_{\alpha,j}$ and random signs for $r_\alpha$.

\subsection{Practical Interpretation}

The large-$K$ theory developed here should be read as the regime in which the
head-specific learning remains order one as $K$ grows \emph{after compensating
for the raw $1/K$ gradient scale of a head-averaged objective}. In that regime,
increasing $K$ turns the ensemble into a deterministic head population whose
law evolves according to \eqref{eq:limit-nu}. This is the regime relevant to
the concrete theorem below and to the population-level interpretation of
ensemble diversity used throughout the note. By contrast, with equal learning
rates across shared and head-specific parameters one has $\lambda_K=1/K\to 0$,
so the leading-order large-$K$ limit freezes the head law instead of transporting
it.

\section{Discussion}
\label{sec:discussion}

The rigorous core of this note is intentionally simple:
\begin{enumerate}[leftmargin=2em]
  \item exact empirical closure at finite $K$;
  \item an abstract bounded-Lipschitz head mean-field theorem;
  \item a concrete finite-horizon large-$K$ theorem for the literal two-layer proxy under bounded-support initialization.
\end{enumerate}

This already covers natural initialization choices in which the shared
parameters are drawn from standard bounded laws such as Xavier/He uniform and
the head-specific parameters are drawn from bounded laws, for instance bounded
output-head initializations for $a_{\alpha,j}$ together with random signs for
$r_\alpha$. In that sense, the bounded-support theorem is not merely a technical toy:
it matches a practically meaningful part of the efficient-ensemble design space.

The message is not that the large-$K$ theory replaces the fixed-$K$,
width-to-infinity collapse picture. Rather, the two limits answer different
questions. The fixed-$K$ width-mean-field theory asks whether a finite set of
heads collapse under deterministic feature learning. The fixed-width,
$K\to\infty$ theory asks whether the head population law remains dispersed or
concentrates.

We do not pursue unbounded or Gaussian head initializations further in the
present draft. Those extensions are interesting, but they substantially
complicate the note and obscure the clean bounded-support story above. Likewise,
we do not try to turn the local heuristic discussion of predictor-law shape,
disagreement, and bias into a global theorem here.

\paragraph{Practitioner-facing reading.}
The regime taxonomy above suggests that ensemble size $K$ should not be viewed
as the only diversity knob. In the large-$K$ picture, increasing $K$ mainly
reduces Monte-Carlo error in approximating the population predictor
\[
  \bar f_t(x)=\int f(x;q_t,z)\,\nu_t(\dd z),
\]
once the underlying law $\nu_t$ is already determined. What changes the regime
itself is the pair consisting of the state law $\nu_t$ and the sensitivity of
the predictor map $H_t:z\mapsto f(\cdot;q_t,z)$ on its support.

This leads to three practical diagnostics:
\begin{enumerate}[leftmargin=2em]
  \item monitor a \emph{disagreement proxy}, such as empirical pairwise
  predictor variance across heads, to distinguish functionally collective from
  functionally diverse behaviour;
  \item compare the ensemble mean to the prediction of the average head state,
  when such an average head is meaningful, to detect bias-relevant regimes;
  \item sweep $K$ jointly with width $m$ or with the strength of head-specific
  parameterization, because a small disagreement signal under large $K$ often
  indicates that the bottleneck is predictor sensitivity rather than sampling
  noise.
\end{enumerate}

In this light, a practitioner who observes little gain from increasing $K$
should first ask whether the ensemble is in a functionally collective regime.
If so, increasing $K$ further is unlikely to help much; stronger head-specific
initialization, earlier head-specific modulation, or a larger shared width $m$
are the more natural interventions. Conversely, if disagreement is appreciable
but the ensemble mean keeps drifting away from the prediction of an average
head, then one is likely in a bias-relevant regime where the head dependence is
genuinely nonlinear.

\paragraph{A simple width heuristic.}
The two-layer proxy also suggests a useful distinction between two possible
sources of head diversity.

\begin{proposition}[Readout-only diversity self-averages with width]
\label{prop:readout-self-average}
Fix a shared state $q=((w_j,b_j)_{j=1}^m)$ and a common modulation vector
$r\in\R^d$. Let
\[
  \phi_j^r(x):=\sigma\bigl(w_j^\top(r\odot x)+b_j\bigr).
\]
For each head $\alpha$, define
\[
  F_\alpha^m(x)
  :=
  \frac1m\sum_{j=1}^m a_{\alpha,j}\phi_j^r(x),
  \qquad
  a_{\alpha,j}=\bar a_j+\xi_{\alpha,j},
\]
where conditional on $(q,r)$ the random variables $\xi_{\alpha,j}$ are
independent across $\alpha$ and $j$, satisfy
\[
  \E[\xi_{\alpha,j}\mid q,r]=0,
  \qquad
  \E[\xi_{\alpha,j}^2\mid q,r]\le v_a,
\]
and let
\[
  \bar F^m(x)
  :=
  \frac1m\sum_{j=1}^m \bar a_j\phi_j^r(x).
\]
Then
\[
  \E[F_\alpha^m(x)\mid q,r]=\bar F^m(x)
  \qquad
  \forall x\in\mathcal X,
\]
and
\begin{align}
  \E\!\left[
    \|F_\alpha^m-\bar F^m\|_{L^2(\mathsf P_X)}^2
    \,\middle|\, q,r
  \right]
  &\le
  \frac{v_a\|\sigma\|_\infty^2}{m},
  \label{eq:readout-only-var}
  \\
  \E\!\left[
    \|F_\alpha^m-F_\beta^m\|_{L^2(\mathsf P_X)}^2
    \,\middle|\, q,r
  \right]
  &\le
  \frac{2v_a\|\sigma\|_\infty^2}{m}
  \qquad
  (\alpha\neq\beta).
  \label{eq:readout-only-pair}
\end{align}
\end{proposition}

\begin{proof}
The conditional mean identity is immediate from
\[
  \E[a_{\alpha,j}\mid q,r]=\bar a_j.
\]
For the variance bound,
\[
  F_\alpha^m(x)-\bar F^m(x)
  =
  \frac1m\sum_{j=1}^m \xi_{\alpha,j}\phi_j^r(x).
\]
Conditional on $(q,r)$, the summands are centered and independent in $j$, so
\[
  \E\!\left[
    |F_\alpha^m(x)-\bar F^m(x)|^2
    \,\middle|\, q,r
  \right]
  =
  \frac1{m^2}\sum_{j=1}^m
  \E[\xi_{\alpha,j}^2\mid q,r]\,
  |\phi_j^r(x)|^2
  \le
  \frac{v_a\|\sigma\|_\infty^2}{m}.
\]
Integrating over $\mathsf P_X$ gives \eqref{eq:readout-only-var}. For
\eqref{eq:readout-only-pair}, note that
\[
  F_\alpha^m-F_\beta^m
  =
  (F_\alpha^m-\bar F^m)-(F_\beta^m-\bar F^m),
\]
and the two centered terms are conditionally independent with the same second
moment bound, so the variances add.
\end{proof}

\begin{corollary}[No Jensen bias from readout dispersion alone]
\label{cor:readout-affine}
Fix $(q,r)$ and define
\[
  H_{q,r}^m(a)(x)
  :=
  \frac1m\sum_{j=1}^m a_j\phi_j^r(x).
\]
Then $a\mapsto H_{q,r}^m(a)$ is affine as a map from $\R^m$ to
$L^2(\mathsf P_X)$. Hence for every probability law $\mu$ on $\R^m$ with finite
first moment,
\[
  \int_{\R^m} H_{q,r}^m(a)\,\mu(\dd a)
  =
  H_{q,r}^m\!\left(
    \int_{\R^m} a\,\mu(\dd a)
  \right).
\]
In particular, dispersion in the output-head coefficients alone cannot create
ensemble bias.
\end{corollary}

\begin{proof}
For fixed $(q,r)$, the map $a\mapsto H_{q,r}^m(a)$ is linear in $a$, so the
claim follows immediately from linearity of Bochner integration.
\end{proof}

These two facts suggest a practitioner-relevant infinite-width heuristic.
Centered head-specific perturbations in the output coefficients $a_{\alpha,j}$
are averaged over the width and therefore tend to generate at most
$O(m^{-1/2})$ predictor differences. By contrast, the modulation vectors
$r_\alpha$ enter \emph{inside} the nonlinear feature map before the averaging
over neurons. If the width-$m$ shared state admits a law-of-large-numbers limit
and the resulting map
\[
  r\mapsto
  \int a\,\sigma\bigl(w^\top(r\odot x)+b\bigr)\,\rho(\dd a,\dd w,\dd b)
\]
is nonconstant on the support of the head law, then one should expect
order-one functional diversity to persist even as $m\to\infty$. Likewise,
nonlinearity of this map in $r$ is the natural source of width-persistent
ensemble bias.

\paragraph{Connection to NTK operating regimes.}
This width heuristic should be read in contrast with the lazy or NTK picture
\citep{embedded-ensembles,chizat2019lazy}. In the NTK regime, parameters remain
close to initialization and the key question is whether the cross-head GP
covariance and off-diagonal NTK vanish. For BatchEnsemble-like architectures,
\citet{embedded-ensembles} identify a special role of the \emph{last-layer
post-activation modulations}: a zero-mean last-layer modulation is sufficient
for model-diagonal GP covariance and NTK, hence for the independent regime in
the infinite-width lazy limit. In that sense, NTK diversity is primarily
\emph{preserved initialization diversity}: the training dynamics propagates the
independence already present in the kernel at $t=0$.

The feature-learning large-$K$ picture is different. By
\cref{prop:mono}, if the initial head law is exactly monomorphic,
\[
  \nu_0=\delta_{z_0},
\]
then the law remains monomorphic for all finite times. Thus deterministic
training does not spontaneously create diversity from a perfectly symmetric
initial condition. What training can do, however, is to \emph{amplify} an
existing state-level spread by changing the predictor map
$H_t:z\mapsto f(\cdot;q_t,z)$. This is precisely why the head-specific
modulation $r$ plays a different role from the output coefficients $a$ in the
present note: centered perturbations in $a$ self-average and remain affine at
the predictor level, whereas perturbations in $r$ enter before the nonlinearity
and can therefore generate width-persistent functional diversity and bias.

This comparison suggests a useful practical distinction between two ways of
obtaining ensemble diversity:
\begin{enumerate}[leftmargin=2em]
  \item \emph{preserved diversity}, where the training dynamics stays close to a
  linearized model and mainly carries forward whatever independence was built
  into the initialization kernel;
  \item \emph{amplified diversity}, where a nontrivial initial head spread is
  transformed by feature learning into stronger function-space separation.
\end{enumerate}
The NTK independent regime is the clean prototype of the first mechanism. The
head-law regime studied here suggests that BatchEnsemble-like models can also
benefit from the second mechanism, but only if the head-specific parameters act
inside feature formation and not merely in the final linear readout.

The most natural next steps are:
\begin{enumerate}[leftmargin=2em]
  \item extend the bounded-support theorem to broader initialization classes;
  \item understand when state-level diversity translates into predictor-level diversity;
  \item characterize disagreement and bias of the large-$K$ predictor in concrete local regimes;
  \item study mixed limits in which both width and ensemble size grow.
\end{enumerate}
\begin{thebibliography}{99}

\bibitem[Blorstad et~al.(2026)Blorstad, Mostein, Blaser, and Parviainen]{blorstad2026}
M.~Blorstad, H.~J.~Mostein, N.~Blaser, and P.~Parviainen.
\newblock Evaluating prediction uncertainty estimates from BatchEnsemble.
\newblock arXiv:2601.21581, 2026.

\bibitem[Chizat and Bach(2018)]{chizat2018}
L.~Chizat and F.~Bach.
\newblock On the global convergence of gradient descent for over-parameterized models using optimal transport.
\newblock In \emph{Advances in Neural Information Processing Systems}, 2018.

\bibitem[Chizat et~al.(2019)Chizat, Oyallon, and Bach]{chizat2019lazy}
L.~Chizat, E.~Oyallon, and F.~Bach.
\newblock On lazy training in differentiable programming.
\newblock In \emph{Advances in Neural Information Processing Systems}, 2019.

\bibitem[Erny(2022)]{erny2022}
X.~Erny.
\newblock Well-posedness and propagation of chaos for McKean--Vlasov equations with jumps and locally Lipschitz coefficients.
\newblock \emph{Stochastic Processes and their Applications}, 150:192--214, 2022.

\bibitem[Gorishniy et~al.(2024)Gorishniy, Kotelnikov, and Babenko]{tabm}
Y.~Gorishniy, A.~Kotelnikov, and A.~Babenko.
\newblock TabM: Advancing tabular deep learning with parameter-efficient ensembling.
\newblock arXiv:2410.24210, 2024.

\bibitem[Gorishniy et~al.(2026)Gorishniy, Kotelnikov, and Babenko]{tabm-repo}
Y.~Gorishniy, A.~Kotelnikov, and A.~Babenko.
\newblock Official TabM repository.
\newblock \url{https://github.com/yandex-research/tabm}, accessed 2026-03-23.

\bibitem[Mei et~al.(2018)Mei, Montanari, and Nguyen]{mei2018}
S.~Mei, A.~Montanari, and P.-M.~Nguyen.
\newblock A mean field view of the landscape of two-layer neural networks.
\newblock \emph{Proceedings of the National Academy of Sciences}, 115(33):E7665--E7671, 2018.

\bibitem[Mei et~al.(2019)Mei, Misiakiewicz, and Montanari]{mei2019}
S.~Mei, T.~Misiakiewicz, and A.~Montanari.
\newblock Mean-field theory of two-layers neural networks: Dimension-free bounds and kernel limit.
\newblock In \emph{Proceedings of the Conference on Learning Theory}, 2019.

\bibitem[Velikanov et~al.(2022)Velikanov, Kail, Anokhin, Vashurin, Panov, Zaytsev, and Yarotsky]{embedded-ensembles}
M.~Velikanov, R.~Kail, I.~Anokhin, R.~Vashurin, M.~Panov, A.~Zaytsev, and D.~Yarotsky.
\newblock Embedded ensembles: Infinite width limit and operating regimes.
\newblock In \emph{Proceedings of the International Conference on Artificial Intelligence and Statistics}, 2022.

\bibitem[Wen et~al.(2020)Wen, Tran, and Ba]{batchensemble}
Y.~Wen, D.~Tran, and J.~Ba.
\newblock BatchEnsemble: An alternative approach to efficient ensemble and lifelong learning.
\newblock In \emph{International Conference on Learning Representations}, 2020.

\bibitem[Zamyatin et~al.(2026)Zamyatin, Indri, Malhotra, and Gartner]{zamyatin2026}
A.~Zamyatin, P.~Indri, S.~Malhotra, and T.~Gartner.
\newblock Is BatchEnsemble a single model? On calibration and diversity of efficient ensembles.
\newblock arXiv:2601.16936, 2026.

\end{thebibliography}

\end{document}
